% TOP_PROBLEM: {"problemId": "problem.separation-of-p-r-from-np-r-in-the-blum-shub-smale-model", "canonicalTitle": "Separation of P_R from NP_R in the Blum–Shub–Smale Model", "releaseRank": 172, "primaryDomain": "applied_computational_mathematics", "primaryDomainLabel": "Applied & computational mathematics", "displayStatus": "open", "releaseStatus": "open"}
% !TeX program = lualatex
% Definition and sources only; this file is not an LLM solution attempt.
\documentclass[11pt]{article}
\usepackage[margin=25mm]{geometry}
\usepackage{fontspec}
\setmainfont{Latin Modern Roman}
\usepackage{amsmath,amssymb,mathtools}
\usepackage{unicode-math}
\setmathfont{Latin Modern Math}
\usepackage{xurl,hyperref}
\hypersetup{colorlinks=true,urlcolor=blue}
\setcounter{secnumdepth}{0}
\setlength{\parindent}{0pt}
\setlength{\parskip}{5pt}
\setlength{\emergencystretch}{3em}
\begin{document}
\section*{\texorpdfstring{Separation of P\_R from NP\_R in the Blum–Shub–Smale Model}{Separation of P\_R from NP\_R in the Blum–Shub–Smale Model}}\label{172}
\subsection{Definitions and mathematical statement}
An ordered-real Blum--Shub--Smale machine stores exact real numbers and performs field operations and equality/order tests at unit cost, using finitely many built-in real constants. Input length is the number of real coordinates, not their binary precision. Let $\mathrm P_{{\mathbb{R}}}$ consist of sets decided in polynomially many deterministic steps. A set belongs to $\mathrm{NP}_{{\mathbb{R}}}$ if membership is equivalent to existence of a polynomial-length real witness accepted by such a verifier. The assertion is
\[
 \mathrm P_{{\mathbb{R}}}\ne\mathrm{NP}_{{\mathbb{R}}}.
\]
The real constants belong to the fixed machine; it does not receive arbitrary input-dependent advice. This is a real-computation separation, not merely the classical binary $\mathrm P$ versus $\mathrm{NP}$ question.
\par\smallskip\textit{Formulation and concept references:} \href{https://www.cs.cmu.edu/~lblum/PAPERS/TuringMeetsNewton.pdf}{[S1]}.

\subsection{Short English statement}
In exact real-number computation, are some problems with short real certificates impossible to solve in polynomially many deterministic operations?
\subsection{Sources}
\begin{itemize}
\item[S1]Computing over the Reals: Where Turing Meets Newton, author exposition associated with2004 publication.
\url{https://www.cs.cmu.edu/~lblum/PAPERS/TuringMeetsNewton.pdf}.
\textit{Formulation source.}
\end{itemize}
\end{document}
